\section{What the policies do, trained and unseen}
\label{sec:results}

This section compares the families on recordings they trained on and recordings
they never saw. Table~\ref{tab:policies} carries the headline numbers,
Figure~\ref{fig:gap} shows the distance between the two halves, and
Figure~\ref{fig:horizon} shows the shape that the single numbers hide. One
property of the comparison has to be settled before any number is quoted, and
it is the subject of the next paragraph.

\paragraph{The families do not plan the same distance ahead.}
Each policy emits a chunk of future commands from one observation, and the
length of that chunk is a property of the family. One emits a hundred steps,
another thirty-two. An error averaged over each policy's own chunk therefore
compares a policy that must commit to more than three seconds of future motion
against one that commits to one second, and reports the difference as skill.
This is not a small correction. Over their own horizons the shorter-horizon
family appears to beat the longer-horizon one by roughly a third; over the
thirty-two steps they share, the margin falls to under a tenth. Every
cross-family claim in this report therefore rests on the shared-horizon figure,
which Table~\ref{tab:policies} reports in its own column beside the native one
so neither can be quoted without the other.

\begin{table}[t]
  \centering
  \caption{Action error against the recorded commands, as a root-mean-square
  deviation. \emph{Horizon} is the number of future steps the family commits to
  from one observation. The left group scores each policy over its own horizon;
  the right group rescores every policy over the shortest horizon any of them
  uses, and is the only group safe to read across rows. \emph{Gap} is held-out
  error divided by trained-on error. It appears in both groups because it
  differs between them: the longer-horizon family's ratio falls from
  \num{5.3} to \num{4.1} when the far end of its chunk --- where its held-out
  error grows fastest --- is excluded.}
  \label{tab:policies}
  \small
  \begin{tabular}{lrrrrrr}
\toprule
& & \multicolumn{3}{c}{over its own horizon} & \multicolumn{2}{c}{over the shared first 32} \\
\cmidrule(lr){3-5}\cmidrule(lr){6-7}
policy & horizon & train & held-out & gap & held-out & gap \\
\midrule
ACT & 100 & 2.672 & 14.091 & 5.3 & 10.107 & 4.1 \\
ACT cropped & 100 & 2.626 & 14.677 & 5.6 & 9.995 & 4.1 \\
Diffusion & 32 & 4.741 & 9.181 & 1.9 & 9.181 & 1.9 \\
Diffusion cropped & 32 & 4.791 & 9.207 & 1.9 & 9.207 & 1.9 \\
pi0.5 & 50 & 17.054 & 18.346 & 1.1 & 14.817 & 1.1 \\
pi0.5 cropped & 50 & 16.968 & 18.108 & 1.1 & 14.621 & 1.1 \\
\bottomrule
\end{tabular}

\end{table}

Read across the shared-horizon column, the two families trained from scratch
differ by less than a tenth, with the diffusion family ahead, and the fine-tuned
family trails both by around half. Read down the gap
column, they differ by more than twofold: the action-chunking family is roughly
four times worse on recordings it never saw than on recordings it trained on,
while the diffusion family is under twice as bad. The gap is a ratio computed
within a single policy, so it is immune to the horizon problem that governs the
rest of the table, and it is the largest and most robust effect reported here.

\begin{figure}[t]
  \centering
  \includegraphics[width=0.85\textwidth]{train_vs_heldout.png}
  \caption{Trained-on against held-out error, every bar scored over the shared
  horizon. The annotation above each pair is the ratio between them. What the
  figure is for is the distance between the two bars of a pair, not the height
  of any bar: the heights depend on the units of the commanded joints, while
  the ratio does not.}
  \label{fig:gap}
\end{figure}

\begin{figure}[t]
  \centering
  \includegraphics[width=\textwidth]{horizon_decay.png}
  \caption{Error against position within the predicted chunk, on a logarithmic
  axis because the families differ by more than an order of magnitude and a
  linear axis collapses two of them onto the floor. The shaded band is the
  shared horizon. The trained-on panel shows the finding that the averaged
  numbers hide: the action-chunking family is almost flat across its entire
  chunk on recordings it trained on, and steep on recordings it did not.}
  \label{fig:horizon}
\end{figure}

Figure~\ref{fig:horizon} shows why a single number per policy is not enough. On
recordings it trained on, the action-chunking family's error barely grows from
the first step of the chunk to the hundredth --- it predicts the far future of a
familiar recording about as accurately as the near future, which is what
reproducing a memorised trajectory looks like and is not what planning looks
like. On recordings it never saw, the same family's error climbs steeply with
distance into the chunk. The diffusion family climbs steeply in both panels,
which is the ordinary and expected behaviour: predictions further ahead are
worse.

\textbf{The fine-tuned family's small gap is not generalisation.} Its held-out
error is only about a tenth above its trained-on error, the smallest gap in
Table~\ref{tab:policies}, and read alone that looks like the most robust policy
in the report. It is not. Over the shared horizon its error on recordings it
\emph{trained on} is \num{13.8}, which is already higher than either other
family's error on recordings it \emph{never saw}, at \num{10.1} and \num{9.2}.
Its two curves in Figure~\ref{fig:horizon} nearly coincide for the same reason.
This family trained for roughly one pass over its recordings, adapting only a
small low-rank addition to a large pretrained model, so it never fitted the
training half closely enough to have anything to lose on the held-out half. Its
row therefore measures what this training budget buys, not how the family
compares at convergence. It also reads three cameras where the others read
five. Its cropped arm is about one per cent lower on both halves, but no
sampler floor was measured for this family, so that difference is not claimed.

\unjustified{} the sampling rate of one frame in twenty was chosen to keep the
scoring of six policies within a working day and has not been shown to leave
the reported errors unchanged. A denser sample would tighten the held-out
figures, whose frame count Section~\ref{sec:recordings} gives as small.
